
\FloatBarrier
\subsection{Image denoising}
\label{sec:denoising}

%\cite{rust.2023}

\ac{SEM} images often contain varying amounts of noise.
While it is preferable to reduce the noise by collecting more signal (i.e. by longer dwell times, higher current, stacking of multiple scans), this is often impossible due to time constraints or the sensitivity of the specimen (see \zcref{fig:ebeam_damage}).

In \ac{SEM} images, multiple noise types can be found.
Shot noise (also known as \textsc{Poisson} noise) arises from the discrete nature of electron detection, leading to random fluctuations in pixel intensity, especially noticeable at low signal levels.
Next to the shot noise originating from the primary and secondary electrons, detector noise can also limit the image quality \cite{agarwal.2024}.
Another typical noise type in the \ac{SEM} is banding noise.
It is caused by the line scanning nature of a \ac{SEM}.
These stripe artefacts in the scanning direction can occur due to object drift or charging.

A relevant pre-processing step for most image processing is therefore denoising.
It can be classified into classical methods, transform techniques and \ac{CNN}-based denoising methods \cite{fan.2019}.
The following paragraphs are intended to provide a brief overview of the most important filters.

Classical methods remove noise in the spatial domain.
Linear filters (e.g. mean filtering) can easily lead to blurring of fine detail and edges in very noisy images.
Non-linear filters like the \textit{median} \cite{huang.1979} and \textit{bilateral} \cite{tomasi.1998} filter provide edge-preserving noise reduction.
\textcite{stutzman.2015} proposed the use of a \textit{thresholded blur} algorithm \cite{schmid.2007}, to preserve edges.
However, this requires the manual selection of a minimum grey value difference to distinguish between phases, which limits its use to images with a significant phase contrast difference.
More advanced classical filters like \ac{NLM} \cite{buades.2005} or total variation \cite{rudin.1992} are better at preserving edges.
However, they also struggle with more complex noise patterns (e.g. low count \ac{EDX}) and can introduce artefacts.

Transform techniques decompose the image into frequency components, allowing for selective noise reduction in specific domains like high-frequency noise (shot noise) or low-frequency, directional patterns (banding noise).
One important representative filter is \ac{BM3D} \cite{dabov.2007}, which is a two-stage filtering method and an extension of the \ac{NLM} filter.
\ac{NLM}-based methods often provide good performance but can be computationally expensive and require a careful selection of parameters.
For this work, the \ac{BM3D} implementation in the extended photo processing module of \textit{OpenCV} \cite{bradski.2000} was used.

\ac{TV} \cite{rudin.1992} filters also deliver good edge preserving results even for low signal images by reducing the total variation of the image signal.
A fast variant of \ac{TV} filters is the split \textsc{Bregman} optimisation-based method \cite{goldstein.2009, getreuer.2012}.
In this work, the implementation provided by the \textit{scikit-image} \cite{vanderwalt.2014} library was used.

More recently, \ac{CNN}-based denoising methods were introduced.
They attempt to learn complex noise patterns directly from degraded-clean image pairs.
These models try to predict the clean image from its noisy counterpart using a training set.
The reduce a loss function and thus, minimise the prediction error.
Various \ac{CNN} architectures have been proposed since 2008 \cite{jain.2008}.
They can be differentiated into multi-layer perceptron models and deep learning models.
The strength of \ac{CNN}s lies in their capacity to learn highly non-linear mappings and adapt to different noise characteristics, often outperforming traditional methods, especially in scenarios with unknown or complex noise distributions.
However, their performance heavily relies on the availability of large and diverse training datasets of degraded-clean image pairs.
In addition, their use is often still a hurdle for untrained users, as they are still rarely available in ready-to-use software.

To objectively evaluate the performance of denoising algorithms, quantitative benchmarks such as \ac{PSNR}, \ac{SSIM} and \ac{MSE} are commonly used.
\ac{PSNR} measures the ratio between the maximum possible signal power and the power of corrupting noise in \unit{\decibel}, providing a simple metric for image fidelity.
\ac{SSIM}, on the other hand, assesses perceived image quality by comparing local patterns of pixel intensities, taking into account luminance, contrast, and texture.
Higher \ac{PSNR} and \ac{SSIM} values indicate better preservation of image quality after denoising.
Finally, \ac{MSE} quantifies the average squared difference between the original and denoised images, with lower values indicating better noise suppression and fidelity to the ground truth.

\textcite{wang.2025} used median, wavelet and a \ac{CNN}-based algorithms (Deep \ac{CNN} with guided filtering) on \ac{SEM}-images of polished sections of hydrated \ac{C3S}.
They came to the conclusion, that the \ac{CNN} significantly outperforms (\ac{PSNR} of \qty{26.35}{\decibel}) the median (\ac{PSNR} of \qty{23.20}{\decibel}) and wavelet filter (\ac{PSNR} of \qty{18.25}{\decibel}).

\FloatBarrier
\subsection{Segmentation of \acs{SEM} and \ac{EDX} images}
\label{sec:lit_segmentation}

Segmentation is the process of partitioning an image into multiple segments (sets of pixels or objects).
The goal is to simplify the representation of the image and make it more meaningful and easier to analyse.
After segmentation, each pixel in the image is assigned to a specific class, cluster or label, which can be used for further analysis.

Segmentation can be done automatically, semi-automatically or manually by a human expert.
While manual segmentation is still used \cite{bajcsy.2020,kleiner.2024a}, it is time-consuming and not reproducible.
The result is highly dependent on the experience and knowledge of the operator, but is often considered as the gold standard to which automated methods are compared to \cite{bajcsy.2020}.

Automated segmentation methods are often used to reduce the time and effort required for image analysis and to enable reproducibility.
However, the prerequisite for a good automatic segmentation is a good image quality.
Nevertheless, especially images acquired using a \ac{SEM} and its detectors are noisy and lack contrast.
These issues can be reduced by applying steps such as denoising, contrast enhancement, and sharpening.
However, low-count \ac{EDX} datasets remain challenging.

Segmentation by applying threshold values to distinguish phases is still a commonly used technique since it is fast and easy to apply \cite{hlobil.2022}.
There are multiple thresholding methods available, such as \textsc{Otsu}'s method \cite{otsu.1979} or \textsc{Kapur}'s method \cite{kapur.1985} to automatically find a threshold to binarise an image.
However, these methods are rarely suitable, since many \ac{SEM} images contain more than two phases or show low contrast.
Additionally, sometimes a specific grey value or range does not represent a specific phase throughout the entire dataset as border regions can appear darker than the image centre.
This is can especially occur in \ac{SEM}-\ac{SE} images, \ac{SEM}-\ac{FIB} or \ac{mCT} datasets.
In those cases, adaptive thresholding methods can be used to find a threshold value for each pixel based on the surrounding pixels \cite{sauvola.2000, phansalkar.2011}.
For example, the method proposed by \textcite{phansalkar.2011} is a good choice for low contrast \ac{SEM}-\ac{SE} images \cite{kleiner.2021a}.
Nevertheless, these methods still require a lot of manual work and subjective decisions to find the right parameters for the specific dataset.

Superpixel segmentation algorithms such as \ac{SLIC} \cite{achanta.2012} divide an image into smaller regions, which are similar in color and texture, while preserving boundaries of objects within the image.
\ac{SLIC}O allows to omit the compactness parameter, since it adaptively chooses this parameter for each superpixel.
These small regions are called superpixels.
These algorithms can help to segment very noisy datasets like from \ac{EDX} \cite{georget.2021, jiang.2022}.
Extensions such as Multi\ac{SLIC} \cite{yu.2023} further enable segmentation across multiple data modalities or data layers resulting in more robust and informative representations.
% TODO: add details from jiang.2022 in the result section

After simplifying the image using superpixel segmentation, the phase assignment must be carried out.
Initial orientation can be achieved through clustering.
Multiple approaches to do this were proposed in recent years. 
Initial orientation can be achieved through clustering using $k$-means.
This can either be done on a per-pixel basis \cite{munch.2015} or on a per-superpixel basis \cite{jiang.2022}.
While not it could not yet applied in the cement context, clustering can be supported by \ac{UMAP} \cite{mcinnes.2018}, which allows to reduce the dimensionality of a multidimensional dataset as found in a multi-element mapping or multiple data sources.
The clustering allows to find similar phase compositions.
The phase assignment for each pixel or superpixel can be done manually or automatically.

using a decision-tree-model, utilising known compositions of different phases \cite{jiang.2022}.

%\cite{stutzman.2004}, vs \cite{georget.2021} vs. unser aachen ansatz \cite{jiang.2022} und vs die blackbox aus aztec \cite{aztec.2004} o.ä software.
% \cite{reed.2025} used a method based on QGIS (see mail from Christiane/Nikkola Paavo, 08.10.2025)

% {
% \color{red}
%Utilizing machine learning methods for segmentation is a promising approach to automate the process and improve the accuracy of the results \cite{}.
% }

\FloatBarrier
\subsection{Sectioning effect}
\label{subsec:sectioning-effect}

Many measurements are made by using \ac{BSE} images of embedded and polished hydrated clinker phases.
In other words, all measurements are taken on random sections through grains and their hydrate layers.
This causes an unavoidable sectioning effect \cite{scrivener.2004}.

In particle size distributions, smaller particles are often over-represented, since the section through the particle tends to occur non-diametrically.
Starting from a monodisperse distribution of spheres with a true diameter of $d_0$, a size distribution with a modal apparent diameter of $d_\text{app} = \pi d_0  / 4$ is obtained \cite{cuzzi.2017}.
% https://doi.org/10.1111/maps.12812 <- 2d to 3d hochrechnen mittels fortran code \cite{cuzzi.2017}
In contrast, the \ac{HLT} around the grain is often overestimated as illustrated in \zcref{fig:sectioning_effect}.

\begin{figure}[htb]
    \centering

	\begin{subfigure}{0.42\textwidth}
		\centerline{\includegraphics[width=.9\textwidth]{hydrate-rim/sectioning_effect_3D.pdf}}
		\caption{
            Two imaging planes (blue) at different sectioning positions within a sphere, where $a \ll b$.
        }
		\label{fig:sectioning_effect_a}

	\end{subfigure}
	\hfill
	\begin{subfigure}{0.57\textwidth}

		\begin{tikzpicture}
			\begin{axis}[
					xmin=0, xmax=1.4,
					ymin=0, ymax=1.4,
					width=0.7\textwidth,
					height=0.7\textwidth,
					grid=major,
					xlabel={\small sectioning distance from centre point $x$},
					ylabel={\small measured radius / \acs*{HLT}},
					ytick={.3,1,1.3},
					xtick={0,.3,.8,1,1.3},
					grid style={line width=.1pt, draw=gray!10},
					legend style={	cells={align=left},
									at={(1.03,.5)},
									anchor=west,
									font=\tiny,
									row sep=4pt},
                    yticklabel style={
                            /pgf/number format/fixed,
                            /pgf/number format/precision=1,
                            /pgf/number format/fixed zerofill,
                    },
                    xticklabel style={
                            /pgf/number format/fixed,
                            /pgf/number format/precision=1,
                            /pgf/number format/fixed zerofill,
                    },
				]

				\addplot[yellow,  domain=0:1, samples=360, very thick] {sqrt(2*1*(x+1) - (x+1)^2)};
				\addlegendentry{grain radius \\ $f(x, r_\text{grain})$};
				\addplot[red, domain=0:1.35, samples=1000, very thick] {sqrt(2*1.3*(x+1.3) - (x+1.3)^2)};
				\addlegendentry{outer radius \\ $f(x, r_\text{outer})$};
				\addplot[blue,   domain=0:1, samples=1000, name path=f] {sqrt(2*1.3*(x+1.3) - (x+1.3)^2) - sqrt(2*1*(x+1) - (x+1)^2)};
				\addlegendentry{\acs*{HLT} \\ $t(x)$};
				\path[name path=axis] (axis cs:0,0) -- (axis cs:1,0);
				\addplot [fill=blue, fill opacity=0.2] fill between[of=f and axis, soft clip={domain=0:1}];

				\node[coordinate, pin={[align=left,pin distance = 5mm]-270:{\tiny $\int_{0}^{r_\text{grain}} t(x) \,dx$}}] at (axis cs:0.4,0.24) {};
				%\addplot[yellow,  domain=0:1, samples=360] {sqrt(2*1*(x+1) - (x+1)^2)};

			\end{axis}
		\end{tikzpicture}
		\caption{
            Influence of the sectioning distance from centre point $x$ on the \acs*{HLT} ($t(x)$), calculated using Equation \zcref{eq:sectioning_effect_2}.
        }
		\label{fig:sectioning_effect_b}
	\end{subfigure}

    \caption{
        Illustration of the sectioning effect on the \acs*{HLT} of a spherical grain.
    }
    \label{fig:sectioning_effect}
\end{figure}

\zcref{eq:sectioning_effect_1, eq:sectioning_effect_2} can be used to calculate the \ac{HLT} ($t$) for a spherical particle, depending on the sectioning distance from the centre point ($x$).

\begin{equation}
	f(x, r) = \sqrt{2r(x+r) - (x+r)^2}
	\label{eq:sectioning_effect_1}
\end{equation}
\begin{conditions}
	f	&	measured radius, \\
	r	&   actual radius, \\
	x	&   sectioning distance from centre point; $0 \leq x \leq r$.
\end{conditions}


% \begin{figure}[htb]
%     \centering
%     \begin{tikzpicture}
%         \begin{axis}[
%             xmin=0, xmax=1.1,
%             ymin=0, ymax=1.1,
%             width=0.7\linewidth,
%             height=6cm,
%             width=6cm,
%             grid=major,
%             xlabel={sectioning distance $x$},
%             ylabel={$f(x,r)/r$ ratio},
%             grid style={line width=.1pt, draw=gray!10},
%         ]
%             \addplot[blue, domain=0:1, samples=100, very thick] {sqrt(1 - x^2)};
% 			\legend
%         \end{axis}
%     \end{tikzpicture}
%     \caption{Ratio of measured radius to actual radius depending on the sectioning distance.}
%     \label{fig:sectioning_effect_ratio}
% \end{figure}

\begin{equation}
	t(x) = f(x, r_\text{outer}) - f(x, r_\text{grain})
	\label{eq:sectioning_effect_2}
\end{equation}
\begin{conditions}
	t				&	measured \ac{HLT}, \\
	r_\text{outer}	&   radius of the grain + hydrate layer, \\
	r_\text{grain}	&   grain radius.
\end{conditions}

\begin{equation}
	M_\text{HLT} = \frac{\int_{0}^{i} t(x) \,dx}{r_\text{grain}}, where~0 < i \leq r_\text{grain}.
	\label{eq:sectioning_effect_3}
\end{equation}
\begin{conditions}
	M_\text{HLT}	&	mean \ac{HLT} of a grain, \\
	r_\text{grain}	&   grain radius, \\
\end{conditions}


\zcref{fig:sectioning_effect_b} indicates that the \ac{HLT} can be significantly overestimated.
This becomes more prevalent the closer the cut through a particle is to the particle border.
% https://www.wolframalpha.com/input?i=integral&assumption=%7B%22F%22%2C+%22Integral%22%2C+%22rangestart%22%7D+-%3E%220%22&assumption=%7B%22C%22%2C+%22integral%22%7D+-%3E+%7B%22Calculator%22%2C+%22dflt%22%7D&assumption=%7B%22F%22%2C+%22Integral%22%2C+%22integrand%22%7D+-%3E%22sqrt%282*1.3*%28x%2B1.3%29+-+%28x%2B1.3%29%5E2%29+-+sqrt%282*1*%28x%2B1%29+-+%28x%2B1%29%5E2%29%22&assumption=%7B%22F%22%2C+%22Integral%22%2C+%22rangeend%22%7D+-%3E%22.9%22
For example, if a grain has a constant \ac{HLT} of \qty{30}{\percent} of the grain radius (as shown in \zcref{fig:sectioning_effect_b}, blue curve), an overestimation of from \qtyrange{0}{23.8}{\percent} of the \ac{HLT} can be expected. %0,3715/0,33
If the grain is cut at a sectioning distance from the centre point of the grain $x$ of \num{0.3}, the grain radius will be underestimated by \qty{4.6}{\percent} of the actual value, while the \ac{HLT} will be overestimated at \qty{1.2}{\percent} of the actual value.
In contrast, at $x=0.8$, the radius will be underestimated by \qty{40.0}{\percent} and the \ac{HLT} will be overestimated at \qty{11.0}{\percent}.

\zcref{fig:hlt_overestimation} illustrates the overestimation of the measured \ac{HLT} depending on the \ac{HLT} / grain radius ratio and $M_\text{HLT}$ (\zcref{eq:sectioning_effect_3}).
Since younger specimens have a smaller \ac{HLT}, a greater overestimation is to be expected, while the overestimation of older specimens with larger layers should be significantly lower.
The dependence on $M_\text{HLT}$ shows that specimens cut at or near the border of the core particle especially increase the error.
This also means that very small particles, which have a higher probability of originating from particle edges, are prone to larger overestimations and should therefore not be considered for further analysis.

\begin{figure}[htb]
    \centering
	\includegraphics[width=.75\linewidth]{hydrate-rim/sectioning_overestimation2.pdf}
	\caption{
        Overestimation of the \ac{HLT} depending on the \ac{HLT} / grain diameter ratio and the integral limit $i$ (see \zcref{eq:sectioning_effect_3}).
        The blue line indicates the results for \ac{HLT} / grain diameter value of \num{0.3}.
    }
	\label{fig:hlt_overestimation}
\end{figure}

\FloatBarrier
\subsection{Calculating the phase composition based on 2D-\acs*{SEM} images}
Using a segmented \ac{SEM} image with the image area $A_1$, the phase composition can be estimated by using the \unit{\areapercent} of alite $A_{\text{C}_3\text{S}}$ and the hydrate phases $A_\text{\text{CSH}}$ and $A_\text{CH}$ \cite{kleiner.2024a}.
If the hydrate phases are indistinguishable, they can be summed up:

\begin{equation}
    \label{eq:est_hydrate_area}
    \begin{aligned}
        A_\text{hydrate} &= A_\text{\text{CSH}} + A_\text{CH} \\
        &= A_1 - A_{\text{C}_3\text{S}} - A_{\text{pores}}
    \end{aligned}
\end{equation}

Although density is a volumetric property and the measurements provide only areas, it can be used to approximate the volume fraction of the hydrate phases.
This approximation assumes that the phase fractions observed in 2D images are representative of the true volumetric fractions.
Under this assumption, Equation \zcref{eq:est_frac} yields an estimate of the mass fraction of the hydrate phases (\ac{CH} and \ac{CSH}), \(\sigma_{\text{hydrate}}\), expressed in \unit{\wtpercent}.

\begin{equation}
    \label{eq:est_frac}
    \sigma_\text{hydrate} = \frac{A_\text{hydrate} \cdot \rho_\text{hydrate}}{A_\text{hydrate} \cdot \rho_\text{hydrate} + A_{\text{C}_3\text{S}} \cdot \rho_{\text{C}_3\text{S}}}
\end{equation}

The densities required to calculate these values are listed in \zcref{tab:densities}.

% why 2.5 : 1.0?
\begin{table}[htb]
    \centering
    \caption{
        Densities of \ac{C3S} \cite{garrault.2006}, \ac{C2S}\cite{zhang.2012a}, \ac{CH} and \ac{CSH} \cite{garrault.2006} different phases in \unit{\gram\per\cubic\centi\meter} \cite{jennings.2000,garrault.2006}.\\$^{*}$mixture of \ac{CSH}:\ac{CH} at a volume ratio of \num{2.5}:\num{1}}
    \label{tab:densities}
    \begin{tabularx}{\textwidth}{@{}XXXXXXX@{}}
        \toprule
        \textbf{Phase} & \ac{C3S} & $\beta$-\ac{C2S} & $\gamma$-\ac{C2S} & \ac{CH} & \ac{CSH} & hydrate phase$^{*}$\\
		\midrule
        \textbf{Density} & \num{3.14} & \num{2.97} & \num{3.28} & \num{2.2} & $\approx$ \num{2.116} &  \num{2.07} \\
		\bottomrule
    \end{tabularx}
\end{table}
% beta-C2S: monoclin, more dense
% gamma-C2S:orthorhombi, less dense

\FloatBarrier
\subsection{Statistical analysis of the representativeness of a 2D image}
\label{sec:2D-representativeness}

A question coming up often, especially in \ac{SEM}, where the field of view is often only a few \unit{\milli\meter\squared} or even less, is how representative an image is of the whole specimen.
However, to evaluate that a measurement is representative can be a rather complicated challenge.
This is especially true for inhomogeneous materials consisting of multiple phases and wide ranges of particle sizes.
The problem starts with sampling and sample preparation.
This process is standardised for granular or powdery samples in multiple fields (e.g. \textcite{din52101.2013} and \textcite{dinen932-1.1996}).
Although this solution is very helpful for analysing loose particles, for the analysis of a monolithic sample the only option is to analyse the largest possible volume or surface area.
This raises the question of what the minimum resolution and the size of the analysed area must be in order to obtain a statistically meaningful result.
\pdfcomment[icon=Comment, color=cyan, author=Christane]{% TODO: comment Christane
	zur statistik von SEM bildanalyse an klinkern gibts auch arbeiten von Paul Stutzman und in dem buch von donald campbell über microscopical examination and interpretation of portland cement clinker bzw die deutsche variante davon. bestimmt gibts auch aussagen von betonmikroskopikern
}

The required lower resolution limit of a \ac{SEM} image can be concluded from the \textcite{nyquist.1928}-\textcite{shannon.1949} theorem. %\textcite{whittaker.1915}-\textcite{nyquist.1928}-\textcite{shannon.1949}
In the context of spatial resolution in digital imaging, this translates to the following requirement:
The pixel should be no larger than half the size of the smallest resolvable feature \cite{roberge.2022}.
In other words, to resolve an object with a size of \qty{4}{\nano\meter}, a resolution of less than \qty{2}{\nano\meter\per\px} is required.
Due to the unavoidable noise level in most digital images (including images acquired by \ac{SEM}, \ac{CT} or light microscopy), a higher resolution is required to reliably represent a feature.
Furthermore, the resolution limitations of the imaging device and the sample preparation (e.g. polishing) and the subsequent segmentation methods have to be considered.

However, the minimum area required to obtain a statistically relevant result is especially an issue for \ac{SEM} analysis, because the imaged field of view is usually only a tiny detail of the whole specimen.
While the operator of the microscope may try to exclude unusual areas or capture a typical, representative area of the specimen, this is a highly subjective task.
Although the operator may try to avoid atypical regions and select a representative area, this judgement is inherently subjective.
Acquiring a larger field of view or several independent images reduces this subjectivity, but it is not always clear how large the dataset must be to attain the desired level of statistical confidence.

In many cases, a particle size distribution or \ac{PSD} is of interest, which allows to evaluate distinct objects of an image \cite{odziomek.2016}.
However, this parameter is not easily applicable for interconnected phases or pore spaces.

Often a reliable phase composition measurement is sufficient.
Therefore, some authors tried to apply a \ac{REA} calculation, to make sure the observed image features are statistically relevant \cite{bosl.1998,kanit.2003,kameda.2006,klaver.2012,houben.2014,jiang.2022}.
Similarly to the \ac{REA}, the size of therepresentative volume element can be calculated.
An essential framework for determining its size was defined by \textcite{kanit.2003}.
%They showed that the required size of the volume element depends directly on the variance of the property under investigation (e.g. phase fraction or porosity) and the desired statistical precision.

In the selected publications to calculate the \ac{REA}, two main methods were used for this statistical analysis (see \zcref{fig:lit_box_methods}).

\begin{figure}[H]
    \centering

	\begin{subfigure}[t]{0.49\textwidth}
		\includegraphics[width=\textwidth]{literature/box_cnt_illust.png}
		\caption{Box count method.}
        \label{fig:lit_box_methods_b}
	\end{subfigure}
	\hfill
	\begin{subfigure}[t]{0.49\textwidth}
		\includegraphics[width=\textwidth]{literature/area growth_illust.png}
		\caption{\ac{BG} method.}
        \label{fig:lit_box_methods_a}
	\end{subfigure}
    \caption{
        Example of the \acs*{BC} and \ac{BG} method applied to a segmented image of \qty{28}{\day} hydrated \ac{C3S}.
        Red: {\color{red}\ac{CSH}}, yellow: {\color{orange}\ac{C3S}}, cyan: {\color{cyan}\ac{CH}}, black: pores.
    }
    \label{fig:lit_box_methods}
\end{figure}

The source code used to compare both methods is publicly available \cite{kleiner.2024b}. %-> https://github.com/kleinerELM/box_count_method

\subsubsection{\Acl*{BC} method}

The \ac{BC} method is a commonly used technique to obtain an indication of the statistical representativeness of the phase composition within an image \cite{mouret.2001, hlavacek.2015, cosenza.2019}.
This method splits an image or multiple images of the same material into smaller sub-tiles (see \zcref{fig:lit_box_methods_b}).
Then, a growing number of tiles is randomly selected from this tile pool.
From this selection, the phase composition is then remeasured.
If this is repeated multiple times, the standard deviation indicates how much area has to be imaged and analysed to obtain a reasonable phase composition measurement.

\subsubsection{\Acl*{BG} method}
\label{sec:box_growth}

A variation of the \ac{BC} method was used by \textcite{kameda.2006} and others \cite{klaver.2012, }. % TODO: at least one more reference
The \ac{BG} method works by randomly selecting a growth point near the centre of an image and gradually increasing the box size around this point.
From this point, a box will grow stepwise as illustrated in \zcref{fig:lit_box_methods_a}.
The phase composition of these growing boxes is continuously measured.
To initiate multiple experiments within an image, the growth point has to be selected randomly within a limited window (one fifth of the image width/height) in the centre of the image.
This allows calculation of a standard deviation.
\textcite{kameda.2006} pointed out, that the data fluctuation of the measured porosity and permeability using this method becomes negligible, when `\textit{the cube size exceeds four grain radii}'.

% relativer Fehler von hlavacek übernehmen

% \cite{cosenza.2019} -> describes REA, box counting method, the statistical method (term from zhang.2000, here called \ac{BG} method)
% \cite{yeong.1998} -> two points correlation function, implemented in Python in: https://github.com/pourion/simulated_annealing/commit/f018edd5e0580d1621826cd3e1afb8f5283a61fa
% variogram (concenza.2019) https://scikit-gstat.readthedocs.io/en/latest/userguide/variogram.html

% lydzba.2014 Microstructure measures and the minimum size of a representative volume element: 2D numerical study
% fauchille.2018

%covariogram
%https://en.wikipedia.org/wiki/Covariance_function
